\label{exp:cech02}

\subsubsection*{Постановка}
\textbf{Физический вопрос.}
Позволяют ли четыре канала кубооктаэдра $(V,\square,E,\triangle)$, дискретная теорема Гаусса $E(k-1\leftrightarrow k) = 12(k^3-(k-1)^3)$ и целочисленные инварианты регистра ячейки ($\Delta\varphi_{\mathrm{disc}}=41$, $N_{\mathrm{HB}}=63$, $\vevHierarchySteps=8$, $N_\varphi=13$) вывести длины всех семи периодов таблицы Менделеева, все семь ядерных магических чисел $\{2,8,20,28,50,82,126\}$, тождество терминального стабильного ядра $^{208}_{82}\text{Pb}_{126}$, нейтронный центр острова стабильности $N_{\mathrm{island}}=184$ и регулярность $1s$-оболочки при $Z=137$ в дискретном графовом операторе Дирака $\mathcal{H}_{\mathrm{Dirac}}^{\mathrm{FCC}}(Z)$ (\cref{thm:chem-electron-and-periodic-table,cor:chem-period8-and-supercritical-limits,thm:chem-fcc-graph-dirac-operator})?

\textbf{Что проверяем.}
\begin{enumerate}[label=\textbf{(\alph*)},leftmargin=*,itemsep=0.25em]
\item Совпадение длин периодов $L_p = 2\lfloor(p+2)/2\rfloor^2 \in \{2,8,8,18,18,32,32\}$, атомных номеров благородных газов $Z_{\mathrm{noble}} \in \{2,10,18,36,54,86,118\}$, ядерных магических полиномов $M_{\mathrm{magic}}^{(\mathrm{low})}(1\ldots 3) = \{2,8,20\}$, $M_{\mathrm{magic}}^{(\mathrm{high})}(4\ldots 7) = \{28,50,82,126\}$, тождества $2\Delta\varphi_{\mathrm{disc}} + 2N_{\mathrm{HB}} = 82 + 126 = 208 = 2\vevHierarchySteps N_\varphi$ ($^{208}_{82}\text{Pb}_{126}$) и $M_{\mathrm{magic}}^{(\mathrm{high})}(8) = 184$.
\item Выполнение дискретного тождества Гаусса $E_{\mathrm{in}}(k) + E_{\mathrm{same}}(k) + E_{\mathrm{out}}(k) \equiv 12(10k^2+2)$ при всех $k\ge 1$ и конечность решёточной гармонической функции Грина в центре: $G_{\mathrm{FCC}}(0) - G_{\mathrm{FCC}}(1) = d = 3$ ($G_{\mathrm{FCC}}(0) \approx 3{,}9158425$).
\item Положительность энергии моды первой сферы $E_{1s}^{(k=1)}(137) > +0{,}6\,m_e c^2$ в графовом операторе Дирака $\mathcal{H}_{\mathrm{Dirac}}^{\mathrm{FCC}}(Z)$ при пробое центральной моды ядра ($k=0$) через $E=0$ при $Z=81\text{--}82$ и через $E=-m_e c^2$ при переходе $Z = 122 \to 123$.
\end{enumerate}

{\sloppy\textbf{Рабочая гипотеза.}
Все перечисленные структурные числа периодической системы, ядерных оболочек и графового спектра Дирака воспроизводятся без свободных параметров (\texttt{Chem\_periodic\_and\_fcc\_dirac}).\par}

\textbf{Наблюдаемые.}
$Z_{\mathrm{noble}}(1\ldots 7)$, $M_{\mathrm{magic}}(1\ldots 8)$, $(Z,N,A)$ для $^{208}_{82}\text{Pb}_{126}$, $G_{\mathrm{FCC}}(0)$, $Z_{\mathrm{core}}(E=0)$, $Z_{\mathrm{core}}(E=-1)$, $E_{1s}^{(k=1)}(137)$ (табл.~\ref{tab:exp-ce-ch-02}).

\textbf{Критерий согласия с теорией.}
Точное совпадение всех целочисленных инвариантов; отсутствие комплексного ветвления при $Z=137$ ($E_{1s}^{(k=1)}(137) > 0{,}5\,m_e c^2$).

\subsubsection*{Теоретическая часть}
По теореме~\ref{thm:chem-fcc-graph-dirac-operator} исключение нижней компоненты переноса~$v_k$ из уравнения $\mathcal{H}_{\mathrm{Dirac}}^{\mathrm{FCC}}(Z)\Psi = E\Psi$ сводит задачу на собственные значения в щели $E \in (-1,+1)$ к невырожденной трёхчленной последовательности Штурма по координационным сферам $k=0,1,\ldots,K$.

\textbf{Проверка размерностей.}
Все атомные и магические числа, решёточная функция Грина $G_{\mathrm{FCC}}(k)$ и отношения энергий $E/(m_e c^2)$ безразмерны:
\[
  \bigl[Z_{\mathrm{noble}}\bigr] = \bigl[M_{\mathrm{magic}}\bigr] = \bigl[G_{\mathrm{FCC}}(k)\bigr] = \left[\frac{E}{m_e c^2}\right] = 1.
\]

\subsubsection*{Ход эксперимента}
{\sloppy Воспроизведение: \texttt{python verify\_principles.py} \texttt{--suite sm} (\texttt{Chem\_periodic\_and\_fcc\_dirac}).\par}

\begin{enumerate}[label=\textbf{Шаг \arabic*:},leftmargin=*,itemsep=0.35em]
\item Вычислить длины периодов, благородные газы, ядерные магические числа $k=1\ldots 8$ и параметры $^{208}_{82}\text{Pb}_{126}$.
\item Проверить тождество рёбер Гаусса на сферах~$S_k$, вычислить $G_{\mathrm{FCC}}(k)$ и решить задачу Штурма для $\mathcal{H}_{\mathrm{Dirac}}^{\mathrm{FCC}}(Z)$.
\item Свести результаты в табл.~\ref{tab:exp-ce-ch-02}.
\end{enumerate}

\begin{table}[htbp]
\centering
\scriptsize
\setlength{\tabcolsep}{2.5pt}
\caption{Верификация периодической системы, ядерных магических оболочек и графового оператора Дирака на \latGCKshort\ (CE-CH-02)}
\label{tab:exp-ce-ch-02}
\begin{tabular}{@{}llrrr@{}}
\toprule
Инвариант / наблюдаемая & Формула ДКМ & Модель & Опыт / эталон & Статус \\
\midrule
Замыкания $Z_{\mathrm{noble}}(1\ldots 7)$ & $\sum_{m=1}^p 2\lfloor(m+2)/2\rfloor^2$ & $(2,10,18,36,54,86,118)$ & $(2,10,18,36,54,86,118)$ & PASS \\
Магические числа $k=1\ldots 7$ & $\frac{k(k+1)(k+2)}{3};\ \frac{k(k^2+5)}{3}$ & $(2,8,20,28,50,82,126)$ & $(2,8,20,28,50,82,126)$ & PASS \\
Ядро $^{208}_{82}\text{Pb}_{126}$ & $2\Delta\varphi_{\mathrm{disc}} + 2N_{\mathrm{HB}} = 2\vevHierarchySteps N_\varphi$ & $82 + 126 = 208$ & $^{208}_{82}\text{Pb}_{126}$ ($T_{1/2}=\infty$) & PASS \\
Остров стабильности ($k=8$) & $M_{\mathrm{magic}}^{(\mathrm{high})}(8) = 8(64+5)/3$ & $N=184,\ Z\in[114,120]$ & $N=184,\ Z\in[114,120]$ & PASS \\
Скачок $G_{\mathrm{FCC}}(0)-G_{\mathrm{FCC}}(1)$ & $d / (1^d - 0^d)$ & $3{,}000000$ & $d = 3$ & PASS \\
Переход $k=0$ через $E=0$ & $E_{\mathrm{core}}^{(k=0)}(81\to 82) < 0$ & $Z=81$ ($E_{82}=-0{,}0455$) & Рубеж $\alpha$-распада $Z>82$ & PASS \\
Пробой $k=0$ в $E=-1$ & $E_{\mathrm{core}}^{(k=0)}(122\to 123) = -1$ & $Z = 122 \to 123$ & Деление $Z\in[120,126]$ & PASS \\
Оболочка $1s$ ($k=1$) при $Z=137$ & Собств.\ значение $\mathcal{H}_{\mathrm{Dirac}}^{\mathrm{FCC}}(137)$ & $+0{,}6331\,m_e c^2$ & Вещественна ($>0$) & PASS \\
\bottomrule
\end{tabular}
\end{table}

\subsubsection*{Анализ, расчёт погрешности и выводы}
Дискретный графовый оператор Дирака на координационных сферах кубооктаэдра полностью свободен от континуальной сингулярности $r\to 0$ при $Z=137$: электронная $1s$-оболочка на первой сфере~$S_1$ ($k=1$) сохраняет положительную энергию $+0{,}6331\,m_e c^2$, тогда как физический предел таблицы Менделеева определяется пробоем центральной моды самого ядра ($k=0$) при $Z=81\text{--}82$ ($E=0$, рубеж свинца $^{208}_{82}\text{Pb}_{126}$) и $Z=122\text{--}123$ ($E=-m_e c^2$, порог спонтанного деления).

\textbf{Вывод.}
Вся периодическая таблица элементов и границы ядерной устойчивости вытекают из геометрии координационных сфер~\latGCKshort.
\textbf{Статус записи:} верифицировано автоматическим гейтом \texttt{Chem\_periodic\_and\_fcc\_dirac}.
